\documentclass{amsart}
% For dealing with references we use the comment environment
\usepackage{verbatim}
\newenvironment{reference}{\comment}{\endcomment}
%\newenvironment{reference}{}{}
\newenvironment{slogan}{\comment}{\endcomment}
\newenvironment{history}{\comment}{\endcomment}

% For commutative diagrams we use Xy-pic
\usepackage[all]{xy}

% We use 2cell for 2-commutative diagrams.
\xyoption{2cell}
\UseAllTwocells

% We use multicol for the list of chapters between chapters
\usepackage{multicol}

% This is generally recommended for better output
\usepackage{lmodern}
\usepackage[T1]{fontenc}

% For cross-file-references

% Package for hypertext links:
\usepackage{hyperref}

% For any local file, say "hello.tex" you want to link to please

% Theorem environments.
%
\theoremstyle{plain}
\newtheorem{theorem}[subsection]{Theorem}
\newtheorem{proposition}[subsection]{Proposition}
\newtheorem{lemma}[subsection]{Lemma}

\theoremstyle{definition}
\newtheorem{definition}[subsection]{Definition}
\newtheorem{example}[subsection]{Example}
\newtheorem{exercise}[subsection]{Exercise}
\newtheorem{situation}[subsection]{Situation}

\theoremstyle{remark}
\newtheorem{remark}[subsection]{Remark}
\newtheorem{remarks}[subsection]{Remarks}

\numberwithin{equation}{subsection}

% Macros
%
\def\lim{\mathop{\mathrm{lim}}\nolimits}
\def\colim{\mathop{\mathrm{colim}}\nolimits}
\def\Spec{\mathop{\mathrm{Spec}}}
\def\Hom{\mathop{\mathrm{Hom}}\nolimits}
\def\Ext{\mathop{\mathrm{Ext}}\nolimits}
\def\SheafHom{\mathop{\mathcal{H}\!\mathit{om}}\nolimits}
\def\SheafExt{\mathop{\mathcal{E}\!\mathit{xt}}\nolimits}
\def\Sch{\mathit{Sch}}
\def\Mor{\mathop{\mathrm{Mor}}\nolimits}
\def\Ob{\mathop{\mathrm{Ob}}\nolimits}
\def\Sh{\mathop{\mathit{Sh}}\nolimits}
\def\NL{\mathop{N\!L}\nolimits}
\def\CH{\mathop{\mathrm{CH}}\nolimits}
\def\proetale{{pro\text{-}\acute{e}tale}}
\def\etale{{\acute{e}tale}}
\def\QCoh{\mathit{QCoh}}
\def\Ker{\mathop{\mathrm{Ker}}}
\def\Im{\mathop{\mathrm{Im}}}
\def\Coker{\mathop{\mathrm{Coker}}}
\def\Coim{\mathop{\mathrm{Coim}}}

% Boxtimes
%
\DeclareMathSymbol{\boxtimes}{\mathbin}{AMSa}{"02}

%
% Macros for moduli stacks/spaces
%
\def\QCohstack{\mathcal{QC}\!\mathit{oh}}
\def\Cohstack{\mathcal{C}\!\mathit{oh}}
\def\Spacesstack{\mathcal{S}\!\mathit{paces}}
\def\Quotfunctor{\mathrm{Quot}}
\def\Hilbfunctor{\mathrm{Hilb}}
\def\Curvesstack{\mathcal{C}\!\mathit{urves}}
\def\Polarizedstack{\mathcal{P}\!\mathit{olarized}}
\def\Complexesstack{\mathcal{C}\!\mathit{omplexes}}
% \Pic is the operator that assigns to X its picard group, usage \Pic(X)
% \Picardstack_{X/B} denotes the Picard stack of X over B
% \Picardfunctor_{X/B} denotes the Picard functor of X over B
\def\Pic{\mathop{\mathrm{Pic}}\nolimits}
\def\Picardstack{\mathcal{P}\!\mathit{ic}}
\def\Picardfunctor{\mathrm{Pic}}
\def\Deformationcategory{\mathcal{D}\!\mathit{ef}}

\setlength{\emergencystretch}{3em}
\begin{document}
\title[Stacks Project, Tag 07EZ]{728 --- Stable cotangent freeness and high powers yield strictly standard presentations}
\maketitle
\noindent Copyright (C) 2005--2025 Johan de Jong. Stacks Project authors. Source: \href{https://stacks.math.columbia.edu/tag/07EZ}{Tag 07EZ}.

\noindent Editorial difficulty note (not author text). Difficulty H2: The author adds generators to stabilize a conormal module, constructs left inverses after localization, and clears denominators to obtain global Jacobian identities. The complete comparison of smooth presentation notions is a specialized algebraic smoothing argument..

\noindent Editorial provenance note (not author text). History: excerpt from revision \texttt{a0444\allowbreak 6e57e\allowbreak c1fbc\allowbreak 252a8\allowbreak 71afc\allowbreak ec775\allowbreak 2fb28\allowbreak 07b14}, smoothing.tex, lines 718--850. Reference presentation was adapted; original-author context and core support blocks were retained or added. The mathematical wording is unchanged. Licensed under GNU FDL 1.2 or later; no invariant sections or cover texts. See ../licenses/COPYING. Context blocks reproduce author definitions and supporting results; they are not additional counted problems.

\begin{definition}
\label{definition-strictly-standard}
Let $R \to A$ be a ring map of finite presentation.
We say an element $a \in A$ is {\it elementary standard in $A$ over $R$}
if there exists a presentation
$A = R[x_1, \ldots, x_n]/(f_1, \ldots, f_m)$
and $0 \leq c \leq \min(n, m)$ such that
\begin{equation}
\label{equation-elementary-standard-one}
a = a' \det(\partial f_j/\partial x_i)_{i, j = 1, \ldots, c}
\end{equation}
for some $a' \in A$ and
\begin{equation}
\label{equation-elementary-standard-two}
a f_{c + j} \in (f_1, \ldots, f_c) + (f_1, \ldots, f_m)^2
\end{equation}
for $j = 1, \ldots, m - c$. We say $a \in A$ is
{\it strictly standard in $A$ over $R$} if there exists a presentation
$A = R[x_1, \ldots, x_n]/(f_1, \ldots, f_m)$
and $0 \leq c \leq \min(n, m)$ such that
\begin{equation}
\label{equation-strictly-standard-one}
a = \sum\nolimits_{I \subset \{1, \ldots, n\},\ |I| = c}
a_I \det(\partial f_j/\partial x_i)_{j = 1, \ldots, c,\ i \in I}
\end{equation}
for some $a_I \in A$ and
\begin{equation}
\label{equation-strictly-standard-two}
a f_{c + j} \in (f_1, \ldots, f_c) + (f_1, \ldots, f_m)^2
\end{equation}
for $j = 1, \ldots, m - c$.
\end{definition}

\begin{definition}
\label{algebra-definition-unramified}
Let $R \to S$ be a ring map.
\begin{enumerate}
\item We say $R \to S$ is {\it unramified} if $R \to S$ is of
finite type and $\Omega_{S/R} = 0$.
\item We say $R \to S$ is {\it G-unramified} if $R \to S$ is of finite
presentation and $\Omega_{S/R} = 0$.
\item Given a prime $\mathfrak q$ of $S$ we say that $S$ is
{\it unramified at $\mathfrak q$} if there exists a
$g \in S$, $g \not \in \mathfrak q$ such that $R \to S_g$ is unramified.
\item Given a prime $\mathfrak q$ of $S$ we say that $S$ is
{\it G-unramified at $\mathfrak q$} if there exists a
$g \in S$, $g \not \in \mathfrak q$ such that $R \to S_g$ is G-unramified.
\end{enumerate}
\end{definition}

\begin{lemma}[Elkik]
\label{lemma-elkik}
Let $R \to A$ be a ring map of finite presentation.
The singular ideal $H_{A/R}$ is the radical of the ideal
generated by strictly standard elements in $A$ over $R$
and also the radical of the ideal generated by elementary
standard elements in $A$ over $R$.
\end{lemma}

\begin{proof}
Assume $a$ is strictly standard in $A$ over $R$. We claim that
$A_a$ is smooth over $R$, which proves that $a \in H_{A/R}$. Namely,
let $A = R[x_1, \ldots, x_n]/(f_1, \ldots, f_m)$, $c$, and $a' \in A$
be as in Definition \ref{definition-strictly-standard}.
Write $I = (f_1, \ldots, f_m)$ so that the naive cotangent
complex of $A$ over $R$ is given by $I/I^2 \to \bigoplus A\text{d}x_i$.
Assumption (\ref{equation-strictly-standard-two})
implies that $(I/I^2)_a$ is generated by the classes of $f_1, \ldots, f_c$.
Assumption (\ref{equation-strictly-standard-one}) implies
that the differential $(I/I^2)_a \to \bigoplus A_a\text{d}x_i$
has a left inverse, see
Lemma \href{https://stacks.math.columbia.edu/tag/07ET}{lemma parse equation strictly standard one (Tag 07ET)}.
Hence $R \to A_a$ is smooth by definition and
Algebra, Lemma \href{https://stacks.math.columbia.edu/tag/00S7}{lemma localize NL (Tag 00S7)}.

\medskip\noindent
Let $H_e, H_s \subset A$ be the radical of the ideal generated by
elementary, resp.\ strictly standard elements of $A$ over $R$.
By definition and what we just proved we have
$H_e \subset H_s \subset H_{A/R}$. The inclusion $H_{A/R} \subset H_e$
follows from Lemma \href{https://stacks.math.columbia.edu/tag/07C6}{lemma find strictly standard (Tag 07C6)}.
\end{proof}


\begin{lemma}
\label{lemma-smooth-standard-smooth}
Let $R \to A$ be a smooth ring map. Then there exists a smooth $R$-algebra
map $A \to B$ with a retraction such that $B$ is standard smooth over
$R$, i.e.,
$$
B \cong R[x_1, \ldots, x_n]/(f_1, \ldots, f_c)
$$
and $\det(\partial f_j/\partial x_i)_{i, j = 1, \ldots, c}$
is invertible in $B$.
\end{lemma}

\begin{proof}
Apply Lemma \href{https://stacks.math.columbia.edu/tag/07CG}{lemma syntomic complete intersection (Tag 07CG)}
to get a smooth $R$-algebra map $A \to C$ with a retraction such that
$C = R[x_1, \ldots, x_n]/(f_1, \ldots, f_c)$
is a relative global complete intersection over $R$. As $C$ is smooth
over $R$ we have a short exact sequence
$$
0 \to
\bigoplus\nolimits_{j = 1, \ldots, c} C f_j \to
\bigoplus\nolimits_{i = 1, \ldots, n} C\text{d}x_i \to
\Omega_{C/R} \to 0
$$
Since $\Omega_{C/R}$ is a projective $C$-module this sequence is split.
Choose a left inverse $t$ to the first map. Say
$t(\text{d}x_i) = \sum c_{ij} f_j$
so that $\sum_i \frac{\partial f_j}{\partial x_i} c_{i\ell} = \delta_{j\ell}$
(Kronecker delta). Let
$$
B' = C[y_1, \ldots, y_c] =
R[x_1, \ldots, x_n, y_1, \ldots, y_c]/(f_1, \ldots, f_c)
$$
The $R$-algebra map $C \to B'$ has a retraction given by mapping $y_j$ to zero.
We claim that the map
$$
R[z_1, \ldots, z_n] \longrightarrow B',\quad
z_i \longmapsto x_i - \sum\nolimits_j c_{ij} y_j
$$
is \'etale at every point in the image of $\Spec(C) \to \Spec(B')$.
In $\Omega_{B'/R[z_1, \ldots, z_n]}$ we have
$$
0 =
\text{d}f_j - \sum\nolimits_i \frac{\partial f_j}{\partial x_i} \text{d}z_i
\equiv
\sum\nolimits_{i, \ell}
\frac{\partial f_j}{\partial x_i} c_{i\ell} \text{d}y_\ell
\equiv
\text{d}y_j \bmod (y_1, \ldots, y_c)\Omega_{B'/R[z_1, \ldots, z_n]}
$$
Since $0 = \text{d}z_i = \text{d}x_i$ modulo
$\sum B'\text{d}y_j + (y_1, \ldots, y_c)\Omega_{B'/R[z_1, \ldots, z_n]}$
we conclude that
$$
\Omega_{B'/R[z_1, \ldots, z_n]}/
(y_1, \ldots, y_c)\Omega_{B'/R[z_1, \ldots, z_n]} = 0.
$$
As $\Omega_{B'/R[z_1, \ldots, z_n]}$ is a finite $B'$-module
by Nakayama's lemma there exists a $g \in 1 + (y_1, \ldots, y_c)$
that $(\Omega_{B'/R[z_1, \ldots, z_n]})_g = 0$. This proves that
$R[z_1, \ldots, z_n] \to B'_g$ is unramified, see
Algebra, Definition \ref{algebra-definition-unramified}.
For any ring map $R \to k$ where $k$ is a field we obtain an
unramified ring map $k[z_1, \ldots, z_n] \to (B'_g) \otimes_R k$
between smooth $k$-algebras of dimension $n$. It follows that
$k[z_1, \ldots, z_n] \to (B'_g) \otimes_R k$ is flat by
Algebra, Lemmas \href{https://stacks.math.columbia.edu/tag/00R4}{lemma CM over regular flat (Tag 00R4)} and
\href{https://stacks.math.columbia.edu/tag/00TS}{lemma characterize smooth kbar (Tag 00TS)}. By the crit\`ere
de platitude par fibre
(Algebra, Lemma \href{https://stacks.math.columbia.edu/tag/00R7}{lemma criterion flatness fibre (Tag 00R7)})
we conclude that $R[z_1, \ldots, z_n] \to B'_g$ is flat.
Finally, Algebra, Lemma \href{https://stacks.math.columbia.edu/tag/00U6}{lemma characterize etale (Tag 00U6)}
implies that $R[z_1, \ldots, z_n] \to B'_g$ is \'etale.
Set $B = B'_g$. Note that $C \to B$ is smooth and has a retraction,
so also $A \to B$ is smooth and has a retraction.
Moreover, $R[z_1, \ldots, z_n] \to B$ is \'etale.
By Algebra, Lemma \href{https://stacks.math.columbia.edu/tag/00U9}{lemma etale standard smooth (Tag 00U9)}
we can write
$$
B = R[z_1, \ldots, z_n, w_1, \ldots, w_m]/(g_1, \ldots, g_m)
$$
with $\det(\partial g_j/\partial w_i)$ invertible in $B$.
This proves the lemma.
\end{proof}


\begin{lemma}
\label{lemma-compare-standard}
Let $R \to A$ be a ring map of finite presentation.
Let $a \in A$. Consider the following conditions on $a$:
\begin{enumerate}
\item $A_a$ is smooth over $R$,
\item $A_a$ is smooth over $R$ and $\Omega_{A_a/R}$ is stably free,
\item $A_a$ is smooth over $R$ and $\Omega_{A_a/R}$ is free,
\item $A_a$ is standard smooth over $R$,
\item $a$ is strictly standard in $A$ over $R$,
\item $a$ is elementary standard in $A$ over $R$.
\end{enumerate}
Then we have
\begin{enumerate}
\item[(a)] (4) $\Rightarrow$ (3) $\Rightarrow$ (2) $\Rightarrow$ (1),
\item[(b)] (6) $\Rightarrow$ (5),
\item[(c)] (6) $\Rightarrow$ (4),
\item[(d)] (5) $\Rightarrow$ (2),
\item[(e)] (2) $\Rightarrow$ the elements $a^e$, $e \geq e_0$ are
strictly standard in $A$ over $R$,
\item[(f)] (4) $\Rightarrow$ the elements $a^e$, $e \geq e_0$ are
elementary standard in $A$ over $R$.
\end{enumerate}
\end{lemma}

\begin{proof}
Part (a) is clear from the definitions and
Algebra, Lemma \href{https://stacks.math.columbia.edu/tag/00T7}{lemma standard smooth (Tag 00T7)}.
Part (b) is clear from Definition \ref{definition-strictly-standard}.

\medskip\noindent
Proof of (c). Choose a presentation
$A = R[x_1, \ldots, x_n]/(f_1, \ldots, f_m)$ such that
(\ref{equation-elementary-standard-one}) and
(\ref{equation-elementary-standard-two}) hold.
Choose $h \in R[x_1, \ldots, x_n]$ mapping to $a$. Then
$$
A_a = R[x_0, x_1, \ldots, x_n]/(x_0h - 1, f_1, \ldots, f_m).
$$
Write $J = (x_0h - 1, f_1, \ldots, f_m)$.
By (\ref{equation-elementary-standard-two}) we see that the $A_a$-module
$J/J^2$ is generated by $x_0h - 1, f_1, \ldots, f_c$
over $A_a$. Hence, as in the proof of Algebra, Lemma \href{https://stacks.math.columbia.edu/tag/07CF}{lemma huber (Tag 07CF)},
we can choose a $g \in 1 + J$ such that
$$
A_a = R[x_0, \ldots, x_n, x_{n + 1}]/
(x_0h - 1, f_1, \ldots, f_c, gx_{n + 1} - 1).
$$
At this point (\ref{equation-elementary-standard-one})
implies that $R \to A_a$ is standard smooth (use the coordinates
$x_0, x_1, \ldots, x_c, x_{n + 1}$ to take derivatives).

\medskip\noindent
Proof of (d). Choose a presentation
$A = R[x_1, \ldots, x_n]/(f_1, \ldots, f_m)$ such that
(\ref{equation-strictly-standard-one}) and
(\ref{equation-strictly-standard-two}) hold.
Write $I = (f_1, \ldots, f_m)$.
We already know that $A_a$ is smooth over $R$, see
Lemma \ref{lemma-elkik}. By
Lemma \href{https://stacks.math.columbia.edu/tag/07ET}{lemma parse equation strictly standard one (Tag 07ET)}
we see that $(I/I^2)_a$ is free on $f_1, \ldots, f_c$
and maps isomorphically to a direct summand of
$\bigoplus A_a \text{d}x_i$. Since
$\Omega_{A_a/R} = (\Omega_{A/R})_a$
is the cokernel of the map
$(I/I^2)_a \to \bigoplus A_a \text{d}x_i$
we conclude that it is stably free.

\medskip\noindent
Proof of (e). Choose a presentation
$A = R[x_1, \ldots, x_n]/I$ with $I$ finitely generated.
By assumption we have a short exact sequence
$$
0 \to (I/I^2)_a \to \bigoplus\nolimits_{i = 1, \ldots, n} A_a\text{d}x_i \to
\Omega_{A_a/R} \to 0
$$
which is split exact. Hence we see that
$(I/I^2)_a \oplus \Omega_{A_a/R}$ is a free $A_a$-module.
Since $\Omega_{A_a/R}$ is stably free we see that $(I/I^2)_a$
is stably free as well. Thus replacing the presentation chosen
above by $A = R[x_1, \ldots, x_n, x_{n + 1}, \ldots, x_{n + r}]/J$ with
$J = (I, x_{n + 1}, \ldots, x_{n + r})$ for some $r$ we get that $(J/J^2)_a$
is (finite) free. Choose $f_1, \ldots, f_c \in J$ which map to a basis of
$(J/J^2)_a$. Extend this to a list of generators
$f_1, \ldots, f_m \in J$. Consider the presentation
$A = R[x_1, \ldots, x_{n + r}]/(f_1, \ldots, f_m)$.
Then (\ref{equation-strictly-standard-two}) holds for $a^e$
for all sufficiently large $e$ by construction. Moreover, since
$(J/J^2)_a \to \bigoplus\nolimits_{i = 1, \ldots, n + r} A_a\text{d}x_i$
is a split injection we can find an $A_a$-linear left inverse.
Writing this left inverse in terms of the basis $f_1, \ldots, f_c$
and clearing denominators we find a linear map
$\psi_0 : A^{\oplus n + r} \to A^{\oplus c}$ such that
$$
A^{\oplus c} \xrightarrow{(f_1, \ldots, f_c)}
J/J^2 \xrightarrow{f \mapsto \text{d}f}
\bigoplus\nolimits_{i = 1, \ldots, n + r} A \text{d}x_i
\xrightarrow{\psi_0}
A^{\oplus c}
$$
is multiplication by $a^{e_0}$ for some $e_0 \geq 1$. By
Lemma \href{https://stacks.math.columbia.edu/tag/07ET}{lemma parse equation strictly standard one (Tag 07ET)}
we see (\ref{equation-strictly-standard-one})
holds for all $a^{ce_0}$ and hence for $a^e$ for all $e$ with $e \geq ce_0$.

\medskip\noindent
Proof of (f). Choose a presentation
$A_a = R[x_1, \ldots, x_n]/(f_1, \ldots, f_c)$ such that
$\det(\partial f_j/\partial x_i)_{i, j = 1, \ldots, c}$
is invertible in $A_a$. We may assume that for some
$m < n$ the classes of the elements $x_1, \ldots, x_m$
correspond $a_i/1$ where $a_1, \ldots, a_m \in A$ are generators of $A$
over $R$, see Lemma \href{https://stacks.math.columbia.edu/tag/07EY}{lemma standard smooth include generators (Tag 07EY)}.
After replacing $x_i$ by $a^Nx_i$ for $m < i \leq n$
we may assume the class of $x_i$ is $a_i/1 \in A_a$ for some
$a_i \in A$. Consider the ring map
$$
\Psi : R[x_1, \ldots, x_n] \longrightarrow A,\quad
x_i \longmapsto a_i.
$$
This is a surjective ring map. By replacing $f_j$ by $a^Nf_j$ we may
assume that $f_j \in R[x_1, \ldots, x_n]$ and that
$\Psi(f_j) = 0$ (since after all $f_j(a_1/1, \ldots, a_n/1) = 0$
in $A_a$). Let $J = \Ker(\Psi)$. Then $A = R[x_1, \ldots, x_n]/J$
is a presentation and $f_1, \ldots, f_c \in J$ are elements such that
$(J/J^2)_a$ is freely generated by $f_1, \ldots, f_c$ and such
that $\det(\partial f_j/\partial x_i)_{i, j = 1, \ldots, c}$
maps to an invertible element of $A_a$. It follows that
(\ref{equation-elementary-standard-one}) and
(\ref{equation-elementary-standard-two})
hold for $a^e$ and all large enough $e$ as desired.
\end{proof}
\end{document}
